\documentclass[12pt,a4paper]{article}
\usepackage{fontspec}
\usepackage{polyglossia}
\setdefaultlanguage{thai}
\setotherlanguage{english}
\setmainfont{THSarabunNew}[Path=/Users/boobee/Library/CloudStorage/OneDrive-UttaraditRajabhatUniversity/Computer Engineering Master Startup - AunQA68/fonts/,Extension=.ttf,UprightFont=*,BoldFont=* Bold,ItalicFont=* Italic,BoldItalicFont=* BoldItalic]
\newfontfamily\thaifont{THSarabunNew}[Path=/Users/boobee/Library/CloudStorage/OneDrive-UttaraditRajabhatUniversity/Computer Engineering Master Startup - AunQA68/fonts/,Extension=.ttf,UprightFont=*,BoldFont=* Bold,ItalicFont=* Italic,BoldItalicFont=* BoldItalic]
\XeTeXlinebreaklocale th
\XeTeXlinebreakskip = 0pt plus 1pt
\usepackage{geometry}\geometry{margin=2cm,top=2cm,bottom=2.5cm}
\usepackage{booktabs,longtable,array,xcolor,multirow,hyperref,colortbl}
\hypersetup{colorlinks=true,linkcolor=blue,urlcolor=blue}
\definecolor{hdr}{RGB}{30,80,140}
\definecolor{pass}{RGB}{200,240,200}
\definecolor{fail}{RGB}{255,215,215}
\definecolor{hi}{RGB}{240,248,255}
\begin{document}
\begin{center}
{\Large\bfseries รายงานผลการประเมิน CLO--PLO Achievement}\\[4pt]
{\large หลักสูตร วศ.ม. วิศวกรรมคอมพิวเตอร์และปัญญาประดิษฐ์}\\[2pt]
{\large มหาวิทยาลัยราชภัฏอุตรดิตถ์ \quad ปีการศึกษา 2568}\\[2pt]
{\normalsize รหัสหลักฐาน: AUNQA-8-4-1}
\end{center}
\vspace{0.2cm}
\noindent\textbf{คำชี้แจงและแหล่งข้อมูลหลักฐาน:} ข้อมูลคะแนนและผลสัมฤทธิ์การประเมินผลลัพธ์การเรียนรู้รายวิชา (CLO) ของนักศึกษารายบุคคลในรายงานฉบับนี้ รวบรวมมาจากรายงานผลการดำเนินการของรายวิชา (มคอ. 5) ทั้งหมด 4 วิชาบังคับของหลักสูตร วศ.ม. วิศวกรรมคอมพิวเตอร์และปัญญาประดิษฐ์ ปีการศึกษา 2568 (ภาคเรียนที่ 1--2) โดยใช้ระดับคะแนนประเมินอิงเกณฑ์ตามเกณฑ์รูบริก (Rubrics Score 4 ระดับ: 1 = ควรปรับปรุง, 2 = พอใช้, 3 = ดี, 4 = ดีเยี่ยม) ที่มีกำหนดไว้ในแผนการสอน (มคอ. 3) ของแต่ละรายวิชา
\vspace{0.3cm}
\section*{1. ความสัมพันธ์ CLO--PLO (ภาพรวมรายวิชา ปีการศึกษา 2568)}
\begin{longtable}{|l|p{6.5cm}|c|c|}
\hline
\rowcolor{hdr}\color{white}\textbf{รหัสวิชา} & \color{white}\textbf{ชื่อวิชา} & \color{white}\textbf{CLO} & \color{white}\textbf{PLO} \\ \hline
\endfirsthead \multicolumn{4}{c}{\tablename\ (ต่อ)} \\ \hline
\rowcolor{hdr}\color{white}\textbf{รหัสวิชา} & \color{white}\textbf{ชื่อวิชา} & \color{white}\textbf{CLO} & \color{white}\textbf{PLO} \\ \hline \endhead
  4095101 & ปัญญาประดิษฐ์และการเรียนรู้ของเครื่อง & CLO1 & PLO1 \\
   &  & CLO2 & PLO1 \\
   &  & CLO3 & PLO3 \\
   &  & CLO4 & PLO3 \\
   &  & CLO5 & PLO1 \\
  \hline
  7015906 & ระเบียบวิธีวิจัยทางวิทยาศาสตร์และวิศวกรรมศาสตร์ & CLO1 & PLO2 \\
   &  & CLO2 & PLO3 \\
   &  & CLO3 & PLO3 \\
   &  & CLO4 & PLO4 \\
   &  & CLO5 & PLO5 \\
  \hline
  4095102 & การเขียนโปรแกรมขั้นสูงสำหรับการเรียนรู้ของเครื่อง & CLO1 & PLO1 \\
   &  & CLO2 & PLO1, PLO2 \\
   &  & CLO3 & PLO2, PLO3 \\
   &  & CLO4 & PLO4 \\
   &  & CLO5 & PLO1, PLO3 \\
  \hline
  7015101 & เทคโนโลยีอุบัติใหม่ทางวิศวกรรมคอมพิวเตอร์และปัญญาประดิษฐ์ & CLO1 & PLO1 \\
   &  & CLO2 & PLO1 \\
   &  & CLO3 & PLO2 \\
   &  & CLO4 & PLO1 \\
   &  & CLO5 & PLO1 \\
  \hline
\end{longtable}
\section*{2. ผลการประเมิน PLO รายนักศึกษา}
\textbf{เกณฑ์:} คะแนน CLO $\geq 60\%$ = ผ่าน และผ่าน $\geq 70\%$ ของ CLO ที่เชื่อมโยง = บรรลุ PLO\\[4pt]
\begin{longtable}{|c|p{6.2cm}|c|c|c|c|}
\hline
\rowcolor{hdr}\color{white}\textbf{PLO} & \color{white}\textbf{คำอธิบาย} & \color{white}\textbf{เกริก} & \color{white}\textbf{พอใจ} & \color{white}\textbf{จาตุร} & \color{white}\textbf{จิรกิ} \\ \hline
\endfirsthead \multicolumn{6}{c}{\tablename\ (ต่อ)} \\ \hline \endhead
  PLO1 & สร้างสรรค์ผลงาน/นวัตกรรมด้วยซอฟต์แวร์/ฮาร์ดแวร์ & \cellcolor{pass}100\% & \cellcolor{pass}100\% & \cellcolor{pass}100\% & \cellcolor{pass}100\% \\ \hline
  PLO2 & ประยุกต์ใช้ความรู้แก้ปัญหาองค์กร/ชุมชน & \cellcolor{pass}100\% & \cellcolor{pass}100\% & \cellcolor{pass}100\% & \cellcolor{pass}100\% \\ \hline
  PLO3 & สร้างองค์ความรู้/ต้นแบบนวัตกรรมด้าน AI & \cellcolor{pass}100\% & \cellcolor{pass}83\% & \cellcolor{pass}100\% & \cellcolor{pass}100\% \\ \hline
  PLO4 & จรรยาบรรณวิชาชีพ & \cellcolor{pass}100\% & \cellcolor{pass}100\% & \cellcolor{pass}100\% & \cellcolor{pass}100\% \\ \hline
  PLO5 & คิดวิเคราะห์/สื่อสาร/ทำงานร่วมกัน/เรียนรู้ตลอดชีวิต & \cellcolor{pass}100\% & \cellcolor{pass}100\% & \cellcolor{pass}100\% & \cellcolor{pass}100\% \\ \hline
\end{longtable}
\section*{3. อัตราการบรรลุ PLO ระดับหลักสูตร (Program-Level Achievement)}
\textbf{เกณฑ์หลักสูตร:} นักศึกษา $\geq 70\%$ ต้องสำเร็จ PLO นั้น\\[4pt]
\begin{longtable}{|c|p{7cm}|c|c|c|}
\hline
\rowcolor{hdr}\color{white}\textbf{PLO} & \color{white}\textbf{คำอธิบาย} & \color{white}\textbf{จำนวน CLO} & \color{white}\textbf{อัตราสำเร็จ} & \color{white}\textbf{ผล} \\ \hline
\endfirsthead \multicolumn{5}{c}{\tablename\ (ต่อ)} \\ \hline \endhead
  \cellcolor{pass}PLO1 & สร้างสรรค์ผลงาน/นวัตกรรมด้วยซอฟต์แวร์/ฮาร์ดแวร์ & 10 & 100\% (4/4) & \cellcolor{pass}\textbf{ผ่าน} \\ \hline
  \cellcolor{pass}PLO2 & ประยุกต์ใช้ความรู้แก้ปัญหาองค์กร/ชุมชน & 4 & 100\% (4/4) & \cellcolor{pass}\textbf{ผ่าน} \\ \hline
  \cellcolor{pass}PLO3 & สร้างองค์ความรู้/ต้นแบบนวัตกรรมด้าน AI & 6 & 100\% (4/4) & \cellcolor{pass}\textbf{ผ่าน} \\ \hline
  \cellcolor{pass}PLO4 & จรรยาบรรณวิชาชีพ & 2 & 100\% (4/4) & \cellcolor{pass}\textbf{ผ่าน} \\ \hline
  \cellcolor{pass}PLO5 & คิดวิเคราะห์/สื่อสาร/ทำงานร่วมกัน/เรียนรู้ตลอดชีวิต & 1 & 100\% (4/4) & \cellcolor{pass}\textbf{ผ่าน} \\ \hline
\end{longtable}
\section*{4. คะแนน CLO รายวิชา — ปีการศึกษา 2568}
\subsection*{4095101 ปัญญาประดิษฐ์และการเรียนรู้ของเครื่อง (ภาค 1/2568)}
\begin{longtable}{|p{4.5cm}|c|c|c|c|c|}
\hline
\rowcolor{hdr}\color{white}\textbf{นักศึกษา} & \color{white}\textbf{CLO1} & \color{white}\textbf{CLO2} & \color{white}\textbf{CLO3} & \color{white}\textbf{CLO4} & \color{white}\textbf{CLO5} \\ \hline
\endfirsthead \hline \endhead
  เกริกเกียรติ บุญไทย & \cellcolor{pass}75\% & \cellcolor{pass}75\% & \cellcolor{pass}75\% & \cellcolor{pass}75\% & \cellcolor{pass}75\% \\ \hline
  พอใจ สุขหา & \cellcolor{pass}75\% & \cellcolor{pass}75\% & \cellcolor{fail}50\% & \cellcolor{pass}75\% & \cellcolor{pass}75\% \\ \hline
  จาตุรงค์ ทรงภักดี & \cellcolor{pass}75\% & \cellcolor{pass}75\% & \cellcolor{pass}75\% & \cellcolor{pass}75\% & \cellcolor{pass}100\% \\ \hline
  จิรกิตติ์ วราภรณ์ & \cellcolor{pass}100\% & \cellcolor{pass}100\% & \cellcolor{pass}100\% & \cellcolor{pass}75\% & \cellcolor{pass}100\% \\ \hline
\end{longtable}
\subsection*{7015906 ระเบียบวิธีวิจัยทางวิทยาศาสตร์และวิศวกรรมศาสตร์ (ภาค 1/2568)}
\begin{longtable}{|p{4.5cm}|c|c|c|c|c|}
\hline
\rowcolor{hdr}\color{white}\textbf{นักศึกษา} & \color{white}\textbf{CLO1} & \color{white}\textbf{CLO2} & \color{white}\textbf{CLO3} & \color{white}\textbf{CLO4} & \color{white}\textbf{CLO5} \\ \hline
\endfirsthead \hline \endhead
  เกริกเกียรติ บุญไทย & \cellcolor{pass}75\% & \cellcolor{pass}75\% & \cellcolor{pass}75\% & \cellcolor{pass}75\% & \cellcolor{pass}75\% \\ \hline
  พอใจ สุขหา & \cellcolor{pass}75\% & \cellcolor{pass}75\% & \cellcolor{pass}75\% & \cellcolor{pass}100\% & \cellcolor{pass}75\% \\ \hline
  จาตุรงค์ ทรงภักดี & \cellcolor{pass}75\% & \cellcolor{pass}75\% & \cellcolor{pass}75\% & \cellcolor{pass}75\% & \cellcolor{pass}75\% \\ \hline
  จิรกิตติ์ วราภรณ์ & \cellcolor{pass}100\% & \cellcolor{pass}100\% & \cellcolor{pass}100\% & \cellcolor{pass}75\% & \cellcolor{pass}100\% \\ \hline
\end{longtable}
\subsection*{4095102 การเขียนโปรแกรมขั้นสูงสำหรับการเรียนรู้ของเครื่อง (ภาค 2/2568)}
\begin{longtable}{|p{4.5cm}|c|c|c|c|c|}
\hline
\rowcolor{hdr}\color{white}\textbf{นักศึกษา} & \color{white}\textbf{CLO1} & \color{white}\textbf{CLO2} & \color{white}\textbf{CLO3} & \color{white}\textbf{CLO4} & \color{white}\textbf{CLO5} \\ \hline
\endfirsthead \hline \endhead
  เกริกเกียรติ บุญไทย & \cellcolor{pass}75\% & \cellcolor{pass}75\% & \cellcolor{pass}75\% & \cellcolor{pass}75\% & \cellcolor{pass}75\% \\ \hline
  พอใจ สุขหา & \cellcolor{pass}75\% & \cellcolor{pass}75\% & \cellcolor{pass}75\% & \cellcolor{pass}100\% & \cellcolor{pass}75\% \\ \hline
  จาตุรงค์ ทรงภักดี & \cellcolor{pass}75\% & \cellcolor{pass}100\% & \cellcolor{pass}75\% & \cellcolor{pass}75\% & \cellcolor{pass}100\% \\ \hline
  จิรกิตติ์ วราภรณ์ & \cellcolor{pass}100\% & \cellcolor{pass}100\% & \cellcolor{pass}100\% & \cellcolor{pass}75\% & \cellcolor{pass}100\% \\ \hline
\end{longtable}
\subsection*{7015101 เทคโนโลยีอุบัติใหม่ทางวิศวกรรมคอมพิวเตอร์และปัญญาประดิษฐ์ (ภาค 2/2568)}
\begin{longtable}{|p{4.5cm}|c|c|c|c|c|}
\hline
\rowcolor{hdr}\color{white}\textbf{นักศึกษา} & \color{white}\textbf{CLO1} & \color{white}\textbf{CLO2} & \color{white}\textbf{CLO3} & \color{white}\textbf{CLO4} & \color{white}\textbf{CLO5} \\ \hline
\endfirsthead \hline \endhead
  เกริกเกียรติ บุญไทย & \cellcolor{pass}75\% & \cellcolor{pass}75\% & \cellcolor{pass}75\% & \cellcolor{pass}75\% & \cellcolor{pass}75\% \\ \hline
  พอใจ สุขหา & \cellcolor{pass}75\% & \cellcolor{pass}75\% & \cellcolor{pass}75\% & \cellcolor{pass}75\% & \cellcolor{pass}75\% \\ \hline
  จาตุรงค์ ทรงภักดี & \cellcolor{pass}75\% & \cellcolor{pass}75\% & \cellcolor{pass}75\% & \cellcolor{pass}100\% & \cellcolor{pass}75\% \\ \hline
  จิรกิตติ์ วราภรณ์ & \cellcolor{pass}100\% & \cellcolor{pass}100\% & \cellcolor{pass}100\% & \cellcolor{pass}75\% & \cellcolor{pass}100\% \\ \hline
\end{longtable}
\vspace{0.3cm}\noindent\rule{\linewidth}{0.4pt}
\noindent\textbf{สรุป:} นักศึกษาที่บรรลุ PLO ครบ 5 รายการ: \textbf{4/4 คน} (100\%)\\
\noindent\textit{ข้อมูล ณ สิ้นปีการศึกษา 2568 (ภาค 1--2) — CLO Aggregation System v1.0 | อัปเดตทุกสิ้นปีการศึกษา}
\end{document}